\documentclass[11pt]{article}
\usepackage[T1]{fontenc}
\usepackage{lmodern,amsmath,amssymb,amsthm}
\usepackage[margin=1in]{geometry}
\usepackage[hidelinks]{hyperref}
\newtheorem{proposition}{Proposition}
\newcommand{\Z}{\mathbb Z}
\newcommand{\Q}{\mathbb Q}
\newcommand{\OO}{\mathcal O}
\newcommand{\LambdaW}{\OO[[W]]}
\setlength{\emergencystretch}{1em}
\title{Ordinary degeneracy maps, period normalization, and a coupled path at conductor 389}
\author{\small\href{mailto:admin@cosmosgeek.org}{admin@cosmosgeek.org}\qquad\href{mailto:b321014@iiit-bh.ac.in}{b321014@iiit-bh.ac.in}}
\date{Working paper for specialist review}
\begin{document}
\maketitle
\begin{abstract}
For the ordinary calculation at $p=13$ on $y^2+y=x^3+x^2-2x$,
we track the integral old quotient after raising the tame level by $7$.
Literal scheme pushforwards introduce the cover degree $6$ in the
period comparison. The two sign matrices give primitive quotient maps
although the old eigensubmodule has nonunit image. We compute the
resulting change of weight and norm coordinates from specified
leading coefficients. Over the old/new nodal algebra an exact
first-order specialization loses the full critical coefficient module
even when a distinguished class has an integral lift. We state the
cohomological inputs to this application separately from the matrix
and power-series identities.
\end{abstract}

\section{The degree in the old period}
Let $q\nmid Np$ be prime, $N\ge4$, and work in weight two on a
trivial tame diamond component. Write $A_1,A_2$ for the two degeneracy
pullbacks from level $N$ to level $Nq$ and $T_1,T_2$ for their
Poincar\'e-adjoint scheme pushforwards, in the conventions of
\cite[\S\S2.3--2.4, 5.7]{KLZ}. Let $a_q$ be the lower Hecke eigenvalue.
The intermediate map that forgets a generator of a cyclic subgroup
of order $q$ has degree $c=q-1$. A fibre of the cyclic-subgroup
correspondence has $q+1$ points, and its two mixed composites give the
usual $q$th Hecke correspondence. Thus on the lower eigenspace the
matrix $(T_i A_j)$ is
\[
 c\begin{pmatrix}q+1&a_q\\a_q&q+1\end{pmatrix}.
\]
Let $\alpha$ be a chosen unit root of $X^2-a_qX+q$ in a
coefficient DVR. Use the compatible Fricke operators $W_H,W_L$,
so $W_H A_1=A_2W_L$, $W_H A_2=A_1W_L$. For
$P=T_1-\alpha^{-1}T_2$ and $P^*=A_1-\alpha^{-1}A_2$ it follows that
\begin{equation}\label{eq:fricke}
 P W_H P^*=c s_q W_L,\qquad
 s_q=\alpha(1-q/\alpha^2)(1-1/\alpha^2).
\end{equation}
Indeed the matrix multiplication gives
$a_q(1+\alpha^{-2})-2(q+1)\alpha^{-1}$
after removing $c$, and $a_q=\alpha+q/\alpha$.
In a normalization where $P^*$ carries the lower normalized
differential to the old stabilized form, the differential functional
satisfies $\omega_L P=\omega_H$. The Fricke-dual functional, with
the same ordinary $p$-root at both levels and normalized by the first
coefficient, then satisfies $\eta_LP=c s_q\eta_H$ by
\eqref{eq:fricke}. In the corresponding determinant frames,
\[
 a_L\det P=c s_q a_H.
\]
This equality specifies which periods are being compared. If raw
pullbacks rather than Poincar\'e-normalized coordinates are used to
identify trivial-diamond parts, the map acquires the relative factor
$c$. Dividing the two pullback identifications by their cover
degrees moves this factor into the Fricke-dual functional. Replacing
$P$ by $P/c$ scales its determinant and the two-fold tensor
pushforward of a diagonal class by the same factor $c^{-2}$.

\section{A finite receiver at the prime 13}
Take $E:y^2+y=x^3+x^2-2x$, $q=7$, and $N=389$. Point counting gives
$a_7=-5$ and $a_{13}=-3$. The selected root of $X^2+5X+7$ is
$\alpha=1158\bmod 2197$, characterized by $\alpha\equiv1\bmod13$.
The following two small blocks are the trace tuples in the recovered
integral modular-symbol calculation at precision $13^3$:
\[
 T^+=\begin{pmatrix}133&348\\32&761\end{pmatrix},\qquad
 T^-=\begin{pmatrix}400&1762\\1203&1795\end{pmatrix}.
\]
Here a row of $T^\pm$ is a literal trace, and the coordinates
are those of the \emph{standard} $U_7$ packet. In this basis
$T_1U_7=7T_2$ and $T_2U_7=a_7T_2-T_1$. Fricke swaps the two traces
and conjugates $U_7$ to the dual Hecke operator $U_7'$. Consequently
$P$ on the Fricke-transported packet is the second row minus
$\alpha^{-1}$ times the first in the displayed coordinates. The results
are
\[
 P^+=(159,829),\qquad P^-=(2064,1710)\pmod {2197}.
\]
Their reductions $(3,10),(10,7)$ are primitive. On the positive
sign the stored matrix and old and new columns are
\[
 U^+=\begin{pmatrix}56&1102\\55&1103\end{pmatrix},\qquad
 o^+=(-1,-1)^t,\qquad n^+=(1418,1)^t .
\]
Direct multiplication gives $P^+U^+=\alpha P^+$, $P^+n^+=0$ and
$P^+o^+=1209\bmod2197$, of valuation one. On the negative sign
$P^-(1896,1)^t=0$. Over a free rank-one old branch quotient a
primitive map to the lower rank-one sign module is an isomorphism.
The extra factor on the old eigen\emph{submodule} reflects the
intersection with the integral new branch.

These packet data, from the original Manin presentation, fix the bases
for this finite calculation. The archived computation uses full plus
ranks $33$ and $262$ at levels $389,2723$, respectively, then
localizes at the non-Eisenstein $T_2=-2$ packet and the positive
$U_7$ root. The sign-minus ambient rank at level $2723$ is $259$.
Both trace-row products and the old column product above were checked
directly from the restored numerical arrays. This checks the arithmetic
of the displayed inputs; the full Manin-basis reduction is from the
original source computation.

\section{A weight derivative and an ordinary annular column}
Use the weight coordinate $W=14^{k-2}-1$ on the level-$389$
ordinary branch. For clarity we describe the finite receiver used
to compute its ordinary root derivative. The arithmetic input here
is the ordinary rank-one control identification and its weight-two
Cartier/Hecke receiver, as in the archived construction using
\cite[\S\S3.3, 3.8]{EPW}. Let $a_n$ be the elliptic newform coefficients
and $C$ the product of $\sum_{n\ge1}\sigma_{11}(n)q^n$ with this newform
(in this paragraph $q$ denotes its Fourier variable).
At weight $14$ the divided $T_{13}$ defect of $E_{12}f$ has
coefficients modulo $13$ given by
\[
 d_n=11(C_{13n}+3C_n)-a_{n/13},
\]
with $a_{n/13}=0$ unless $13\mid n$. To see it, use
$E_{12}\equiv1+13\cdot11\sum\sigma_{11}(n)q^n\bmod169$.
The weight-$14$ return term $13^{13}$ is zero modulo $169$; at
weight two it is $13$, giving the final term above. Cartier squared
minus the identity gives coefficients $d_{169n}-d_n$ in the
source receiver. Its $U_{13}$ residue is $10$, so its
Cartier-difference eigenvalue $10^2-1=8\bmod13$ is inverted by $5$.

Here are finite instructions for checking the scalar. The archived
coefficient of $T_n$ in the projector, for $1\le n\le30$ in order,
is
\[
\begin{gathered}
 (2,1,9,9,5,2,0,8,5,3,11,3,8,12,7,4,7,0,4,9,7,3,8,2,3,1,8,5,5,2).
\end{gathered}
\]
Form the $a_n$ through $65910$ by point counts, the good prime-power
recurrence and multiplicativity (the multiplicative coefficient at
$389$ is $1$). Sieve $\sigma_{11}(n)\bmod13$ through the same bound,
and convolve at indices needed in the displayed defects. For the
projector coefficients $e_n$ the finite sum
$\sum e_n(d_{169n}-d_n)$ gives $12\bmod13$; this convolution and
projection were recomputed using the recovered coefficient vector.
The $T_{13}$ secant divided by $13$ is thus $5\cdot12=8\bmod13$.
At weight $14$ the nonordinary root is zero modulo $169$ and the
ordinary root is $101\bmod169$. At weight two the ordinary root of
$X^2+3X+13$ equals $114\bmod169$. Since $W(14)=156\bmod169$,
the corresponding derivative is
\[
 \alpha_{13}'(0)=
 \frac{(101-114)/13}{156/13}=1\pmod{13}.
\]
In the convention $\mathcal L_{\rm GS}=-2\,d\log\alpha_{13}/dk$
this yields $\mathcal L_{\rm GS}/\log_{13}(14)=5\bmod13$.
The finite checks do not use the complex central derivative.

Another coefficient comes from an annular theta-series test. At the
inert prime write its integral unimodular local norm form as
$Q(z_1,z_2)=\mathrm N(z_1)+u\mathrm N(z_2)$ with $u\in\mathbb Z_p^\times$
and both coordinates in the integers of the unramified quadratic field.
On fibres $Q\in p\mathbb Z_p$, either both coordinates are units or
both are divisible by $p$. Thus the annular condition that the first
coordinate be a unit equals, on these fibres, the difference between
the full lattice and the sublattice scaled by $p$. If the non-$p$
test functions are preserved by this change of variables, the theta
series for the scaled sublattice is $V_{p^2}\Theta_0$. Since $U_p$
selects exactly the $p$-divisible fibres, passage to a common level
and ordinary projection give
\[
 e_{\rm ord}\Theta_{\rm ann}=(1-U_p^{-2})e_{\rm ord}\Theta_0.
\]
Indeed $U_p^2V_{p^2}=1$ and $U_p$ is invertible on the ordinary
summand. The multiplier on the retained residual packet equals
$1-10^{-2}=11\bmod13$. The calculation requires the same
Schwartz function, tame phase and lattice frame on both sides.

\section{The old/new congruence quotient}
Let $\Lambda=\LambdaW$ and normalize the weight determinant by
$\delta(W)$, with $\delta(0)=1$. Write $\nu=\delta^{1/2}$,
$\alpha^2-a_q(W)\alpha+q\delta=0$, and $d(W)=\alpha/\nu-1$.
With the two standard old bases, consider
\[
 U_i=\begin{pmatrix}0&-\delta\\q&a_q\end{pmatrix},\qquad
 U_o=\begin{pmatrix}a_q&-\delta\\q&0\end{pmatrix},\qquad
 G=\begin{pmatrix}q+1&a_q\\a_q/\delta&q+1\end{pmatrix}.
\]
The vectors $v_i=(-\delta,\alpha)^t$, $v_o=(\alpha,q)^t$ are
$\alpha$-eigenvectors, and
$Gv_i=\alpha^{-1}(\alpha^2-\delta)v_o$. This last scalar is
$h d(W)$ for $h=\nu(\alpha+\nu)/\alpha$, a unit in the positive
residual packet at odd $p$.

Here is the lattice consequence and its necessary hypotheses.
Let $M$ be a free ordinary sign module over $\Lambda$, and let
$i:\Lambda\to M$, $j:M\to\Lambda$ be the old inclusion and the dual
trace in the selected root. Assume $i(1)$ is residually primitive,
$j$ surjective, and $ji=hd$. These are integral Ihara inputs
(see \cite[\S2.2]{Wiles} and \cite[\S3]{EPW} for the degeneracy and
ordinary control frameworks). Suppose on the fraction space
$T=U_q/\nu-1$ acts by $d$ on $i\Lambda$, by zero on $\ker j$, and
$d\ne0$. Both are direct summands over the fraction field. Thus
$T=i h^{-1}j$ there, and consequently over the original lattice
by torsion-freeness. It follows that
\[
 TM=i\Lambda,qquad \ker T=\ker j,qquad
 M/(TM+\ker T)\simeq\Lambda/(d).
\]
This argument allows a new summand of arbitrary rank. The coefficient
algebra of a two-branch rank-two node is
$B_W=\Lambda[t]/(t(t-d(W)))$; multiplication by $t$ in this algebra
has $\operatorname{ann}(t)=(t-d)B_W$. Thus a free $B_W$-module of
rank two gives two copies of $\Lambda/(d)$ in the analogous quotient.

Let $d(0)\equiv0\bmod\pi^n$ and $D=\OO/(\pi^n)$.
Under $W\mapsto\eta$ in $A=D[\eta]/(\eta^2)$, write $a=d'(0)\bmod\pi^n$
and $v=\min(n,v_\pi(a))$. For $Q=\Lambda/(d)$ the free two-term
resolution yields
\[
 Q\otimes A=A/(a\eta),\qquad
 \operatorname{Tor}^{\Lambda}_1(Q,A)
     =\operatorname{ann}_{D}(a)\oplus\eta D .
\]
The map $Q\otimes D\to\eta(Q\otimes A)$ is $D\to D/(a)$.
The kernel has length $n-v$ and the Tor group length $n+v$.
Changing the generator of $(d)$ by a unit leaves both lengths invariant.
For a unit $d'(0)$ the entire critical quotient dies under multiplication
by $\eta$. This loss occurs on the congruence quotient before any
comparison to a unitary Hecke algebra.

\section{A diagonal norm factor and the coupled path}
We give the source calculation in a form that separates its
input identities from its consequences. Let $s=\sigma_q^{-1}$ denote
the normalized tame Frobenius variable of a diagonal Rankin class.
In the raw degeneracy normalization the four terms from the
Beilinson--Flach norm relation \cite[\S5]{KLZ} on a scalar
diagonal weight-two packet have factors
\[
 1-s^2,qquad a_qs(1-s),\qquad
 a_qs(1-s),\qquad qs(1-s^2).
\]
Contracting both factors by $P$ gives
\[
 (1-s^2)-2\alpha^{-1}a_qs(1-s)+q\alpha^{-2}s(1-s^2)
   =(1-s)^2(1-q\alpha^{-2}s).
\]
For $s=(1+Y)^{-\kappa_q}$, $\kappa_q=\log_p\langle q\rangle/
\log_p(1+p)$, the initial $Y^2$ coefficient is
$\kappa_q^2(1-q\alpha^{-2})$. At $q=7$, $p=13$ one has
$\kappa_7=4\bmod13$, giving $8\bmod13$. For instance
$\kappa_q\bmod13$ is
$(q^{12}-1)/(13\cdot12)\bmod13$. A bottom class of order two
in $Y$ is thereby sent to order four if the coefficient on the lower
determinant line is a unit and the quotient projection is an integral
isomorphism in the specified frames.

For the coupled weight/norm construction take $\OO=\mathbb Z_{13}$.
Assume the retained normalized diagonal source on $B_W$ has the
following form, where the indicated Tate line is free on $e$:
\begin{equation}\label{eq:nodal}
 z=Y^3 C(W,Y)(t-d(W)+X(W,Y))e,\qquad
 C\in\OO[[W,Y]]^\times,quad X\in Y\OO[[W,Y]]^\times .
\end{equation}
Here the last condition means $X/Y$ is a unit. The weight character
in the determinant of the rank-two module is canceled by choosing
the external variable $1+Y_j=(1+W)(1+Y)$. The data on both branches
and in the same full node in \eqref{eq:nodal} are required. Knowing
only two vanishing orders at $W=0$ does not establish this identity.

\begin{proposition}
If $d_0=d(0)\ne0$, $d_0\in13\OO$ and $d'(0)$ is a unit, there is
a unique integral power series $W(Y)\in Y\OO[[Y]]$ with
$d(W(Y))-X(W(Y),Y)=d_0$. On the new branch $t=0$ the quotient of
$z$ by $d_0Y^3$ along this path is primitive. If the normalized old,
new and lower level determinant frames satisfy
\[
 a_{\rm new}/a_H\equiv-1\bmod13,qquad
 X_Y(0,0)\equiv-32\frac{d_0 a_{\rm new}}{a_L\det P}\pmod{13},
\]
and $d_0=\alpha-1$, $d'(0)=5\bmod13$ as in the recorded $389,7$
normalization, then $X_Y(0,0)=1$ and $W'(0)=8$ modulo $13$.
\end{proposition}
\begin{proof}
The implicit equation has unit derivative in $W$ since $X(W,0)=0$.
Solving coefficient by coefficient gives existence and uniqueness.
Its substitution in \eqref{eq:nodal} gives $-d_0Y^3 C(W(Y),Y)e$.
The frame congruence in the proposition arises at a rank-two node as
follows. In its coefficient dual use $\tau$ with $\tau(1)=0$ and
$\tau(t)=1$. The old and new normalized lines are $t\tau,(t-d_0)\tau$.
Resolving $\tau$ along these two lines gives coefficients $1/d_0$
and $-1/d_0$. A further common unit of the node takes values congruent
modulo $d_0$ on the branches. Thus the determinant ratio has sign
$-1$ modulo $13$ when the old and new frames differ only by this common
unit. The scalar $-32$ incorporates the normalized source coefficients
of the lower and new classes; it is an input from the archived
reciprocity calculation (see the scope statement below).

From \eqref{eq:fricke} we have $a_L\det P=c s_q a_H$, and
$s_q/d_0=(1-q/\alpha^2)(\alpha+1)/\alpha\equiv1\pmod{13}$.
Thus $d_0a_{\rm new}/(a_L\det P)\equiv-1/6=2\bmod13$ and
$X_Y(0,0)\equiv-32\cdot2=1$. Implicit differentiation gives $W'(0)=1/5=8$.
\end{proof}

Here is the distinction between a class lift, a Hecke character and
a module lift. For $h=h_0+t h_1$ in the node,
$h\mapsto h_0+X h_1$ respects the defining polynomial exactly when
$X(X-d)=0$ in the target ring. Its first-order multiplicativity
defect is $-d_0 X_Y(0,0)\eta$ at fixed weight. A first-order
character therefore survives after the $d_0$ cut. Also
$(t-d)^2=-d(t-d)$: passing from a primitive new eigenline to the
determinant of the rank-two new eigen\emph{submodule} costs $d$.
Along the coupled path put $S=\OO[[Y]]$, $f=d_0+X(W(Y),Y)$ and
$Q_{\rm path}=S/(f)$. Weierstrass preparation gives $Q_{\rm path}\simeq\OO$,
with $Y$ acting by a nonzero scalar of valuation $v_{13}(d_0)$.
On reducing $d_0=0$, $Y^2=0$, the equation becomes a unit times $Y$:
the scalar quotient is $\OO/d_0$ and the analogous rank-two quotient
is $(\OO/d_0)^2$. This full coefficient loss is consistent with a
primitive lift of the one class in \eqref{eq:nodal}.

In fact, for $V=S/(d_0)$ one has
\[
 \mathrm{Hom}_S(Q_{\rm path},V)=\mathrm{Hom}_S(V,Q_{\rm path})=0,qquad
 \operatorname{Ext}^1_S(Q_{\rm path},V)\simeq
 \operatorname{Ext}^1_S(V,Q_{\rm path})\simeq\OO/d_0.
\]
This follows by the principal resolutions: $f\bmod d_0$ is $Y$ times
a unit, and $d_0$ is regular on $Q_{\rm path}$. The critical fibres
agree, even though their identity need not lift to a map of the modules.
For a lift of the form $v(Y)=-C(Y)\kappa(Y)$ in a free source,
where $\kappa=\kappa_0+Y\kappa_1+\cdots$ is a cocycle and the
cochain differential is $d+Y\theta+\cdots$, we instead have
$d\kappa_1+\theta\kappa_0=0$. The corresponding equation for
$v_1=-C_0\kappa_1-C_1\kappa_0$ gives the explicit first-order
nullhomotopy for that class.

\section{The opposite local factor and the scope of the calculation}
The old split place here is $q_0=389$ over $K=\Q(\sqrt{-11})$.
If $J$ denotes the independent norm coordinate, alignment of the
outer character requires $(1+W)(1+J)^2=1+Y$. Thus
$J'(0)=(1-W'(0))/2=3\bmod13$. In the local spectral coordinate this
fixes $y=q_0^{-1}(1+Y)^{\kappa_{q_0}}$. For the normalized
intertwiner on its common kernel, the scalar of the opposite
composition in this coordinate is the rational function
\[
 r(y)=\frac{(1-y/q_0^2)(1-y/q_0^3)}{1-q_0^{-3}/y}.
\]
This scalar, the domain being the common kernel before any finite cut,
is a local input. Its behavior along the path can be computed directly.
Here $d_0=143\bmod169$. The small roots of $d(W)$ and $f(Y)$
satisfy $W_0=39$, $\rho=26\bmod169$, respectively, by the linear
terms since $W_0,\rho\in13\OO$. With $C=(1+\rho)^{\kappa_{389}}$,
$\kappa_{389}=9\bmod13$, one finds $C=66\bmod169$. At this point
\[
 (1-Cq_0^{-4})/13=5,\quad (1-q_0^{-2}/C)/13=10,\quad
 1-Cq_0^{-3}=2\quad(\bmod 13).
\]
Hence $r(q_0^{-1}C)=1\bmod13$, while at $Y=0$ the simplified
$r(q_0^{-1})=(1-q_0^{-3})(1+q_0^{-2})$ has residue $4$.
A unit value at one point gives neither an everywhere integral family
nor the opposite input itself. In the archived normalization, if $f_0$
is the coherent critical vector and $D(q_0)$ the integral common-kernel
operator, the opposite vector involves
$(1-q_0^{-3})^{-1}(1-q_0^{-1})^{-1}D(q_0)f_0$, as well as its scalar.
It is not obtained by simply identifying the two vectors.

\paragraph{What can be reproduced.}
The original November 17 record preserves the two sign projection
matrices, two independent $U_{13}$ arithmetic outputs, the source
Fourier coefficients, the tame norm polynomial, and the original
node/Tor computations. Their small matrix products, convolution, slope
and residue calculations are given above. The broader source's
identification of the control projector, the normalized lower and
new adjoint/regulator inputs giving $-32$, and the all-weight class
factorization \eqref{eq:nodal} are separate inputs to the stated path
application. A later reconstructed check of the period units does not
retain every intermediate raw calculation; no independent reconstruction
of the complete tame Cartan numerator or the whole Beilinson--Flach
class is asserted here. The finite calculations involve $p=13$ on
the level-$389$ source raised at $7$. The conductor-$389$ higher
moment with coefficient prime $389$ is a different calculation.
Finally no original unitary Hecke character or identification of its
central class follows from the native old/new normalization, the
existence of a source class or its first-order nullhomotopy.

\begin{thebibliography}{9}
\bibitem{KLZ} G. Kings, D. Loeffler and S. L. Zerbes,
\emph{Rankin--Eisenstein classes and explicit reciprocity laws},
Cambridge J. Math. \textbf{5} (2017), 1--122.
\href{https://arxiv.org/abs/1503.02888}{Corrected version}.
\bibitem{Wiles} A. Wiles, \emph{Modular elliptic curves and Fermat's
Last Theorem}, Ann. of Math. \textbf{141} (1995), 443--551.
\href{https://doi.org/10.2307/2118559}{Article}.
\bibitem{EPW} M. Emerton, R. Pollack and T. Weston,
\emph{Variation of Iwasawa invariants in Hida families},
Invent. Math. \textbf{163} (2006), 523--580.
\href{https://www.math.uchicago.edu/~emerton/pdffiles/iwasawa.pdf}{Authors' paper}.
\end{thebibliography}
\end{document}
